\chapter{Size-Invariant Sparsity --- and Why the Plan's Proposed Fix Fails (Phase 2.2)}
\label{chap:p2-2}

\section{Goal}

The pilot's visual-sparsity metric measures concentration as the fraction of attention mass held by the
top-5 of 576 attention cells. Chapter~6 of the pilot report found this measure carries a strong size
confound: the correlation between $\log(\mathrm{box\ area})$ and the top-5 mass ratio is $-0.70$, so a
large object (an excavator or a guardrail spanning much of the frame) is scored as ``unfocused''
largely \emph{because} it is big --- its attention necessarily spreads over more than five cells. The
scale-up plan (item 4.2) proposes a fix: judge the attention mass \emph{inside the grounded box}
against \emph{the box's own cell footprint} rather than against all 576 cells, so an object covering
40\% of the frame is judged against 40\% of the cells.

The goal of this phase was to implement that fix and confirm it removes the confound. It does not ---
and establishing \emph{that}, then finding a measure that does, is the phase's real contribution.

\section{What Was Done}

Three candidate concentration measures were implemented (each a pure, unit-tested function) and their
size confound measured directly as the correlation with $\log(\mathrm{box\ area})$ over the 156 scored
pilot samples:

\begin{enumerate}
  \item \texttt{box\_mass\_ratio} --- the plan's exact proposal: mass inside the box $\div$ the box's
  footprint fraction. Under \emph{uniform} attention this is $1.0$ at any box size (the property its
  unit tests confirm).
  \item \texttt{box\_mass\_inside} --- the mass inside the box \emph{without} the footprint division:
  the fraction of total attention that lands on the object, counting all its cells.
  \item \texttt{gini\_concentration} --- the Gini coefficient of the whole attention map, a measure of
  the \emph{shape} of the attention distribution that does not reference the object box at all.
\end{enumerate}

All three live behind script~06's \texttt{-{}-size-invariant} flag, which writes a suffixed CSV so the
frozen \texttt{visual\_sparsity.csv} is preserved.

\section{Results}

\subsection{The plan's formula over-corrects; mass-inside flips the bias}

\begin{center}
\begin{tabular}{p{5.6cm}p{3.2cm}p{3.4cm}}
\toprule
\textbf{measure} & \textbf{corr with $\log$ area} & \textbf{verdict} \\
\midrule
top-5 mass ratio (pilot)          & $-0.755$ & reproduces the $-0.70$ size bias \\
mass $\div$ footprint (plan)      & $-0.828$ & \emph{over}-corrects (worse) \\
mass inside box                   & $+0.819$ & flips the bias the other way \\
Gini of the map (box-independent) & $\mathbf{-0.437}$ & confound roughly halved \\
\bottomrule
\end{tabular}
\captionof{table}{Size confound of each concentration measure, as its correlation with
$\log(\mathrm{box\ area})$. The pilot metric's $-0.755$ reproduces the reported $-0.70$. The plan's proposed
mass-over-footprint ratio makes the confound \emph{stronger}, not weaker; only the box-independent Gini
substantially reduces it.}
\label{tab:p22-confound}
\end{center}

The mechanism is clear once measured. Dividing by the box footprint over-corrects because real
attention is peaked, not uniform: a small object's box covers only a cell or two, so a tiny footprint
in the denominator inflates the ratio, and the measure ends up \emph{more} anti-correlated with size
($-0.83$) than the metric it was meant to fix. Removing the division instead (mass inside) flips the
sign: a large box simply captures more of the map's mass, so big objects now score \emph{high}
($+0.82$). Neither is size-invariant; both are size-\emph{dominated}, in opposite directions. The
lesson is that no box-relative mass scalar can be size-invariant here, because object size and
attention spread are entangled through a fixed $24\times24$ grid.

\subsection{The measure that actually works does not use the box}

The Gini coefficient of the attention map describes how unequal the attention is across cells --- pure
distribution shape, computed without any reference to the object's box or size. Its correlation with
$\log$ area is $-0.437$: not zero, but a substantial reduction from the pilot metric's $-0.755$. More
tellingly, its per-rule pattern no longer tracks size:

\begin{center}
\begin{tabular}{p{3.2cm}p{2.6cm}p{2.6cm}p{2.8cm}}
\toprule
\textbf{rule} & \textbf{top-5 ratio} & \textbf{Gini} & \textbf{typical box} \\
\midrule
rule\_1 (hard hat)  & 0.125 & 0.603 & small \\
rule\_2 (harness)   & 0.113 & 0.576 & small \\
rule\_3 (guardrail) & 0.098 & 0.629 & \textbf{large} \\
rule\_4 (excavator) & 0.079 & 0.553 & \textbf{large} \\
\bottomrule
\end{tabular}
\captionof{table}{Per-rule concentration. Under the pilot's top-5 measure the largest-object rule
(rule\_3) scores \emph{lowest} (0.098) --- the artifact. Under Gini it scores \emph{highest} (0.629),
and the four rules sit in a narrow 0.55--0.63 band that no longer orders by object size.}
\label{tab:p22-perrule}
\end{center}

Under the pilot metric the largest-object rule looked the most unfocused; under Gini it is the most
focused, and the size ordering is gone. The residual $-0.44$ correlation is real but modest, and is
better read as genuine focus signal (with a small remaining size effect) than as the artifact the
$-0.70$ represented --- a large object's attention does spread slightly more, lowering its Gini a
little, but not enough to recreate the ``big means unfocused'' inversion.

\section{The Problem, and the Resolution}

The surfaced problem is that a plan item, taken literally, does not achieve its stated goal --- the
mass-over-footprint ratio was a reasonable hypothesis that the data falsifies. The resolution was to
test it rather than trust it: implement the proposed formula, measure its confound, watch it fail,
and then reason about \emph{why} (footprint division over-corrects on peaked attention) to arrive at a
measure that sidesteps the entanglement entirely by not using the box. Gini is adopted as the
size-robust sparsity measure; the two box-relative measures are retained in the output for reference
and transparency, explicitly labelled as size-biased so no future reader mistakes either for the fix.
Where a single headline sparsity number is quoted, it should additionally be reported within
object-size bands, since even Gini's residual $-0.44$ means size has not vanished entirely.

\section{Standard vs. Non-Standard Choices}

\textbf{Standard.} Using a distribution-inequality statistic (Gini) as a concentration measure is
ordinary; so is checking a metric for a confound by correlating it with the suspected nuisance
variable.

\textbf{Non-standard, and the contribution.} The phase treats its own specification as a hypothesis to
be tested, not a change to be implemented. The plan's proposed size-invariant formula is shown
empirically to make the confound worse, and is replaced only after a measured alternative is validated
against the same $\log$-area correlation the pilot used --- with the failed candidates kept in the
record, labelled, rather than quietly dropped. Falsifying a plausible-sounding fix and documenting the
mechanism of its failure is worth more to the eventual paper than silently shipping the formula would
have been.

\section{Implications}

Visual sparsity now has a size-robust measure (Gini) whose per-rule comparisons are no longer an
artifact of object size, so cross-rule focus claims can be made without the $-0.70$ caveat that shadowed
the pilot's. The finding also sharpens a general principle for the remaining Phase~2 metrics: any
box-IoU or box-mass quantity is a candidate for the same size entanglement, which is exactly why the
plan calls for size controls in stability (2.3) and robustness (2.4) as well. Gini, being
box-independent, is additionally available for the grid-fallback samples that have no grounded box at
all, where the box-relative measures cannot be computed.
